\documentclass[11pt,a4paper]{article}
\usepackage[utf8]{inputenc}
\usepackage{amsmath,amssymb,amsfonts}
\usepackage{geometry}
\usepackage{hyperref}

\geometry{margin=1in}

\title{\textbf{\Huge Directed Graph Feedback Loop Stability \& Memory Buffers}\\[0.5em]
\Large Real-Time DSP Safety Guarantees in Dynamic Modular Graphs}
\author{\textbf{Time Dilation DAW (v0.0.1) Engineering Group}}
\date{\today}

\begin{document}

\maketitle

\begin{abstract}
This document formulates the directed graph cycle detection algorithms, topological sorting methods, and 1-block delay buffer history allocation mechanism ensuring $100\%$ real-time audio thread stability in \textbf{Time Dilation DAW}.
\end{abstract}

\section{Modular Audio Directed Graph Representation}
The modular audio processing network is modeled as a Directed Graph $G = (V, E)$, where $V$ represents audio and control processing nodes and $E \subseteq V \times V$ represents patch cable connections:
\begin{equation}
E = \{ (u, v, i_{\text{out}}, j_{\text{in}}) \mid \text{Output } i_{\text{out}} \text{ of Node } u \text{ connects to Input } j_{\text{in}} \text{ of Node } v \}
\end{equation}

\section{Topological Sort \& DFS Feedback Loop Detection}
A feedback loop occurs when there exists a directed cycle $v_1 \to v_2 \to \dots \to v_k \to v_1$ in $G$.
Depth-First Search (DFS) with node visitation states $\text{State}(v) \in \{\text{UNVISITED}, \text{VISITING}, \text{VISITED}\}$ identifies cyclic dependencies in linear $O(|V| + |E|)$ time complexity:
\begin{equation}
\text{IsFeedback}(u \to v) = \text{true} \iff \text{State}(v) == \text{VISITING}
\end{equation}

\section{1-Block History Delay Buffer Allocation}
For any connection $e = (u, v) \in E$ identified as part of a feedback loop cycle:
\begin{equation}
x_{\text{dest}}[n] = y_{\text{src, prevBlock}}[n], \quad \forall n \in [0, K_{\text{block}}-1]
\end{equation}
where $y_{\text{src, prevBlock}}$ is the audio buffer rendered during block $N-1$.
For non-cyclic (causal forward) connections:
\begin{equation}
x_{\text{dest}}[n] = y_{\text{src, currentBlock}}[n], \quad \forall n \in [0, K_{\text{block}}-1]
\end{equation}

\section{Mathematical Proof of Thread Safety \& Zero Starvation}
Let $K_{\text{block}}$ be the audio buffer block size (e.g., 512 samples).
Because all cyclic dependencies read from $y_{\text{src, prevBlock}}$, the input buffer contents $x_{\text{dest}}$ for cyclic edges are completely determined prior to processing block $N$.
Therefore, the dependency graph $G_{\text{acyclic}} = (V, E \setminus E_{\text{feedback}})$ forms a Directed Acyclic Graph (DAG), guaranteeing:
\begin{enumerate}
\item \textbf{Deterministic Topological Evaluation}: Node execution order is strictly well-ordered.
\item \textbf{Zero Audio Thread Deadlocks}: No processing node waits for uncalculated live samples within the same buffer quantum.
\item \textbf{Absolute Memory Boundedness}: Memory overhead per feedback edge is strictly limited to $1 \times K_{\text{block}}$ float array allocation.
\end{enumerate}

\end{document}
